\documentclass{article}
\usepackage{graphicx} % Required for inserting images

% <added stuff>
\usepackage{amssymb}
\usepackage{amsmath}

\newcommand{\R}{\mathbb{R}}
\newcommand{\N}{\mathbb{N}}
\newcommand{\Z}{\mathbb{Z}}
\newcommand{\Q}{\mathbb{Q}}
\newcommand{\C}{\mathbb{C}}

\newcommand{\power}{\mathcal{P}}
\newcommand{\T}{\mathcal{T}}
\newcommand{\F}{\mathcal{F}}
\newcommand{\G}{\mathcal{G}}
\newcommand{\B}{\mathcal{B}}
\newcommand{\U}{\mathcal{U}}
\newcommand{\mcO}{\mathcal{O}}

\newcommand{\eps}{\varepsilon}
\newcommand{\nullset}{\varnothing}
\newcommand{\lb}{\left(}
\newcommand{\rb}{\right)}
\newcommand{\minus}{\backslash}
\newcommand{\supp}{\text{supp}}
% </added stuff>

\title{Fundamentals 28}
\author{Siddharth K}
\date{May 2025}

\begin{document}

\maketitle

\subsection*{Problem 28.1}
Let $\{a, b, c\} = X$. Since there are only three possible values for the size of the smallest element in a filter, there are three filters on $X$ up to relabeling:
\begin{enumerate}
    \item $\{X\}$
    \item $\{\{a, b\}, X\}$
    \item $\{\{a\}, \{a, b\}, X\}$
\end{enumerate}

\subsection*{Problem 28.2}
We prove the three axioms of a filter:
\begin{enumerate}
    \item $\B$ only contains nonempty subsets of $X$ so $\nullset \notin \F$.
    
    \item Let $F_1, F_2 \in \F$. We know that $F_1 \supseteq B_1$ and $F_2 \supseteq B_2$ where $B_1, B_2 \in \B$, so $F_1 \cap F_2 \supseteq B_1 \cap B_2 \supseteq B_3$ for some $B_3 \in \B$. Since $F_1 \cap F_2$ is a superset of an element of $\B$, $F_1 \cap F_2 \in \F$.
    
    \item Since $\F$ is generated by $\B$, it is the set of all supersets of $\B$, so every $F \in \F$ has $G \in \F$ where $F \subseteq G$.
\end{enumerate}

\subsection*{Problem 28.3}
\begin{enumerate}
    \item $\nullset \not\supseteq A$, so $\nullset \notin \F$.
    \item Let $F, G \in \F$. Then $F \cap G \supseteq A$, so $F \cap G \in \F$.
    \item Let $F \in \F$ and $F \subseteq G$. Then $G \supseteq F \supseteq A$ so $G \in \F$.
\end{enumerate}

\subsection*{Problem 28.5}
Let $B_1 = f(F_1)$ and $B_2 = f(F_2)$ for any $F_1, F_2 \in \F$. We know that $F_1 \cap F_2 \in \F$, so $f(F_1 \cap F_2) \in \B$. We have that $f(F_1 \cap F_2) \subseteq f(F_1) \cap f(F_2) = B_1 \cap B_2$. So we can let $B_3 = f(F_1 \cap F_2) \in \B$, so that $B_3 \subseteq B_1 \cap B_2$.

\subsection*{Problem 28.6}
Assume there is finer $\G \supseteq \F$ with $\G \to x$. Since 
$F \subseteq G$, every open neighborhood of $x$ is in $\G$ and therefore $\F$ as well. So $x$ is a limit point of $\F$. \\

Now assume that $x$ is a limit point of $\F$, so that $x \in \bigcap_{F \in \F} \overline{F}$. Let $\B = \{U \cap F : F \in \F, \;U \text{ open neighborhood of } x\}$. \\

Let $U$ be an open neighborhood of $x$. Then $U \supseteq U \cap F \in \F$, so $U \in \G$. Also, $\G \supseteq \F$ because, if we let $F \in \F$, then $U \cap F \subseteq F$ and $U \cap F \in \G$ so $F \in \G$. Therefore $G \to x$.

\subsection*{Problem 28.8}
Let $\F$ be an ultrafilter and $\G \supseteq \F$. For contradiction, assume there's $G \in \G \minus \F$. Since $G \notin \F$, $X \minus G \in \F$, so $X \minus G \in \G$. But then $\G$ contains $G$ and $X \minus G$, which is impossible since $\G$ is a filter. So we have a contradiction, and $\F$ is the finest filter on $X$.

\end{document}
